\section{Introduction}

Congressional redistricting combines a fixed numerical obligation---create
exactly $k$ districts---with a vast collection of geographically feasible
assignments. Recursive bisection reduces that choice space to a sequence of
two-way decisions. For California's 52-district delegation, the canonical
schedule begins
\[
52\rightarrow 26/26,\qquad 26\rightarrow 13/13,\qquad
13\rightarrow 6/7.
\]
The schedule is intelligible, reproducible, and naturally parallel. The
remaining discretion lies in selecting each cut.

Practical graph partitioners such as METIS are well suited to finding strong
balanced cuts quickly \cite{karypis1998fast}. Their output, however, is not an
optimality proof. Re-running a deterministic heuristic can establish what the
program produced; it cannot establish that no better cut was permitted by the
same rules.

This paper asks whether the recursive structure itself can become
proof-carrying. Our answer is conditional but constructive. A governing rule
must freeze the unit universe, population, adjacency and island links, split
schedule, objective order, and tie-break. Software can then certify each
required cut and bind the cuts into one recursive tree. This follows the
procedural lesson of Huntington--Hill: mathematics settles execution only
after the governing criterion has been selected \cite{huntington1928}.

The contribution is a methods and evidence contract. We do not claim that the
chosen geographic objective is inevitable, that the resulting plan is a legal
safe harbor, or that certified plans already outperform all alternatives on
compactness, representation, or runtime.
